%!TEX root = paper.tex

\section{Inseparability Relations}

We start by giving an alternative characterization of $\S$-query and $\S$-fact inseparability 
that will allow us to prove our results in a more uniform and clear way.

\begin{proposition} \label{prop:fq-insep-charact-rules}
TBoxes $\T$ and $\T'$ are
\begin{enumerate}
\item $\S$-query inseparable iff $\T\models \rrule \Leftrightarrow \T'\models \rrule$ holds for every rule $\rrule$ over $\S$;
\item $\S$-fact inseparable iff $\T\models \rrule \Leftrightarrow \T'\models \rrule$ holds for every datalog rule $\rrule$ over $\S$.
\end{enumerate}
\end{proposition}
\begin{proof}
 It suffices to observe that, for every TBox $\T$ and every rule $\rrule=\fml\to\fmm$ over $\S$, 
 $\T\models\rrule$ iff $\T\cup\mset{\fct\s}{\fct\in\fml}\models\fmm\s$, with $\s$ a substitution 
 mapping all free variables in $\rrule$ to fresh, pairwise distinct constants.
\end{proof}

\begin{proposition} \label{prop:insep-rel-hierarchy}
${\eqm}\subsetneq{\eqq}\subsetneq{\eqf}\subsetneq{\eqi}$.
\end{proposition}
\begin{proof}
  The inclusion ${\eqq}\subseteq{\eqf}$ is immediate by definition
  while ${\eqf}\subseteq{\eqi}$ follows by
  Proposition~\ref{prop:fq-insep-charact-rules}. The inclusion
  ${\eqm}\subseteq{\eqq}$ follows since
  $\ax^{\calI}=\ax^{\calI|_{\S}}=\ax^{\calI'|_{\S}}=\ax^{\calI'}$ for
  every rule $\ax$ over $\S$ whenever $\calI$ and $\calI'$ coincide on
  $\S$.

  To show strictness of the inclusions, we can w.l.o.g. restrict
  ourselves to the signature $\S=\set{\roleQ,\bot}$ where $\roleQ$ is
  a unary predicate (if $\S$ contains more symbols, one can consider
  $\T$ such that $\sig{\T}\cap\S\subseteq\set{\roleQ}$; adapting the
  argument to higher arities for $\roleQ$ is also straightforward;
  finally, the presence of $\bot$ in $\S$ is not relevant for the
  proof).

  For ${\eqm}\subsetneq{\eqq}$, suppose $\T=\set{\top(x)\to\exists
  y.[\roleR(x,y)\land\A(y)],\top(x)\to\exists y.[\roleR(x,y)\land\B(y)],\A(x)\land\B(x)\to\roleQ(x)}$. 
  Then $\T\eqa\emptyset$.
  However, $\T\not\eqm\emptyset$ since, for any interpretation $\calI$ with a singleton domain 
  such that $\roleQ^{\calI} = \emptyset$, $\calI$ cannot be turned into a model of $\T$ 
  without changing the interpretation of $\roleQ$.

  For ${\eqq}\subsetneq{\eqf}$, suppose $\T=\set{\top(x)\to\exists y.[\roleR(x,y)\land\roleQ(y)]}$. 
  Then $\T\eqf\emptyset$ but $\T\not\eqq\emptyset$ since $\T\models\exists x.\roleQ(x)$ 
  while $\emptyset\not\models\exists x.\roleQ(x)$.

%  For ${\eqa}\subsetneq{\eqr}$, suppose $\T=\set{\top(x)\to\exists y.[\roleR(x,y)\land\roleQ(y)]}$. 
%  Then $\T\eqr\emptyset$ but $\T\not\eqa\emptyset$ since $\T\models\exists x.\roleQ(x)$ 
%  while $\emptyset\not\models\exists x.\roleQ(x)$.
%
%  For ${\eqr}\subsetneq{\eqd}$, suppose $\T=\set{r}$ where
%  $r=\roleQ(a)\to\roleQ(b)\lor\roleQ(c)$. Then $\T\eqd\emptyset$
%  but $\T\not\eqr\emptyset$ since $\T\models r$ while $\emptyset\not\models r$.

  For ${\eqf}\subsetneq{\eqi}$, suppose $\T=\set{r}$ where
  $r=\roleQ(a)\land\roleQ(b)\to\roleQ(c)$. Then $\T\eqi\emptyset$ but
  $\T\not\eqd\emptyset$ since $\T\models r$ while $\emptyset\not\models r$.
  \end{proof}











% \section{Embedding Hyperresolution Proofs in Forward Chaining Proofs}
\section{Deductive Inseparability}


%This section contains the arguments for correctness of Theorems 

Theorems \ref{thm:datalogAtomicSubsumptionModules}, \ref{thm:disjDatalogModules},
\ref{thm:datalogCQModules},
\ref{thm:datalogClassifModules}
are all shown by a similar argument, which we present next.

\subsubsection{Hyperresolution}
Given $\ax = \fml(\x)\to \exists \y. [\bigvee_{i=1}^n
\fmm(\x,\y)]\in\T$\/ we denote with $\sk(\ax)$ the result of applying
standard Skolemisation to $\ax$---which replaces, for each $y\in\y$,
all occurrences of $y$ in $\ax$ by $f_y(\x)$, where $f_y$ is a fresh
function symbol unique for $y$.  Given a substitution $\theta$ mapping
existentially quantified variables in $\T$ to constants and a
Skolemised formula $\fml$, we write $\Gamma_\theta(\fml)$ for the
formula obtained from $\fml$ by replacing every occurrence of a
functional term $f_y(\t)$ by the constant $y\theta$.
%  we denote with 
% $\Gamma_{\theta}$ the mapping that transforms any first order formula $\fml$ into the result of replacing all 
% occurrences in $\fml$ of functional terms of the form $f_{y}(\t)$ by the constant $y\theta$.

By distributing disjunctions over conjunctions in the head of
$\sk(\ax)$ we obtain a rule of the form $\fml \to \bigwedge_{j=1}^m
\fmm'_j$ where each $\fmm'_j$ is a disjunction of atoms.  We denote
with $\cnf(\ax)$ the set $\mset{\fml \to \fmm'_j}{1\leq j\leq m}$ and
extend this notation in the natural way to $\cnf(\T) =
\bigcup_{\ax\in\T}\cnf(\ax)$.  We call $\cnf(\T)$ a \emph{CNF TBox}
and each $\srule\in\cnf(\T)$ a \emph{CNF rule}.
%$\T$ and $\cnf(\T)$ are equisatisfiable; 
%furthermore, 
Clearly, $\cnf(\T)\models\T$, and hence $\T\cup\abox\models \fml'$
implies $\cnf(\T)\cup\abox\models \fml'$ for every $\abox$ and
$\fml'$.


%\begin{definition}
Let $\fml$ be a disjunction of facts, $\abox$ an ABox,
and $\srule=\bigwedge_{i=1}^n\fct'_i\to\fmm\in\cnf(\T)$.
A formula $\fml$ is a \emph{hyperresolvent} of $\srule$ and ground disjunctions 
\mbox{$\fct_1\vee \fmm_1,\dots,\fct_n\vee \fmm_n$} (with each $\fmm_i$ potentially empty)
if $\fml = \bigvee_{i=1}^n\fmm_i \vee \fmm\s$  with $\s$
a MGU of $\fct_i,\fct'_i$ for each $1 \leq i \leq n$.
Let $\cT$ be a CNF TBox. A hyperresolution proof (or simply a \emph{proof}) 
of $\fml$ in $\cT\cup \abox$ is a pair $\rho = (T, \lambda)$ where $T$ is a tree, 
$\lambda$ is a mapping from nodes in $T$ to disjunctions of  facts,
and from edges in $T$ to CNF rules in $\cT$, such that 
for every node $v$ the following properties hold:
\begin{enumerate}
\item $\lambda(v) = \fml$ if $v$ is the root of $T$;
\item $\lambda(v) \in \abox$ if $v$ is a leaf in $T$;
and 
\item if $v$ has children
$w_1,\dots,w_n$ then each edge from $v$
to $w_i$ is labelled by the same CNF rule $\srule$
and  $\lambda(v)$ is a hyperresolvent of $\srule$
and $\lambda(w_1), \dots, \lambda(w_n)$.
\end{enumerate}
If there exists a proof of $\fml$ in $\cT\cup\abox$ we write $\cT\cup\abox\vdash\fml$.
The \emph{support} of $\fml$ is the set of CNF rules occurring in some proof of $\fml$ in $\cT\cup\abox$.
%Let $\rel(\fml, \cnf(\T), \abox)$ denote the set of rules that support some proof
%of $\fml$ in $\cnf(\T)\cup\abox$. We extend this notation naturally to forward-chaining proofs, 
%defined in Section \ref{sec:prelim}: for a fact $\fct$ and a datalog program $\prog$, 
%$\rel(\fct, \prog, \abox)$ will denote the set of rules that support some proof of $\fct$ in 
%$\prog\cup\abox$.

Hyperresolution is sound (if $\cT\cup\abox\vdash\fml$ then $\cT\cup\abox\models\fml$)
and complete in the following sense:
if $\cT\cup\abox\models \fml$ then there exists $\fmm\subseteq \fml$ such that 
$\cT\cup\abox \vdash \fmm$.
%\end{definition}

\vspace{4mm}

Given a \bla{} $\upchi$ and $\rrule\in\T$, we denote with $\Xi^{\chi}(\rrule)$ the set of datalog rules in 
$\prog^{\chi}$ corresponding to $\rrule$, as described in Definition \ref{def:bla}.
The following auxiliary results provide the basis for 
correctness of our approach to module extraction.


\begin{restatable}{lemma}{hyperresolutionProofsToDatalogProofs}
\label{lemma:HyperresolutionProofsToDatalogProofs}
Let $\upchi = \langle \theta, \abox_0, \arel \rangle$ be a \bla{}.
Let $\N$ be the set of constants mentioned in $\upchi$. 
Let $\abox$ be a function-free ABox that only mentions constants that are
fresh w.r.t. $\T$ and $\N$.
Let $\nu$ be a mapping from constants in $\abox$ to $\N$ such that 
$\abox\nu\subseteq \abox_0$.
Let $\fml$ be a disjunction of facts and $\rho = (T, \lambda)$ a proof of 
$\fml$ in $\cnf(\T)\cup\abox$. 
The following properties hold:
\begin{enumerate}
\item $\prog^{\chi}\cup\abox_0 \vdash \Gamma_{\theta}(\fct\nu)$ for every $\fct\in\fml$.
\item For every $\ax\in\T$ such that $\rho$ mentions some $\srule\in \cnf(\ax)$
there exists $\fct\in \fml\cup\{\bot\}$ and a proof of $\Gamma_{\theta}(\fct\nu)$ 
in $\prog^{\chi}\cup\abox_0$ that mentions some rule in $\Xi^{\chi}(\ax)$. 
\end{enumerate}
\end{restatable}
%\hyperresolutionProofsToDatalogProofs*
\begin{proof}
We reason by induction on the depth $d$ of $\rho$.
\begin{itemize}
\item [$d=0$]\ \\
In this case $\fml$ must be a fact in $\abox$. Since $\abox$ is function-free by assumption 
we have $\Gamma_{\theta}(\fml\nu) = \fml\nu$, and since $\abox\nu\subseteq \abox_0$
we have $\fml\nu\in\abox_0$. Therefore, there exists a trivial proof of $\Gamma_{\theta}(\fml\nu)$ 
in $\prog^{\chi}\cup\abox_0$ and property $1$ is satisfied. Furthermore, if the depth 
of $\rho$ is $0$ then there cannot be any rules in its support, so property $2$ is 
trivially satisfied as well. 

\item [$d>0$]\ \\
Let $v$ be the root of $T$ and $w_1,\dots,w_n$ the children of $v$. 
Then it must be
\begin{itemize}
\item $\lambda(w_i) = \fcu_i \vee \fmm_i$ for each $1\leq i \leq n$;
\item $\lambda(v,w_i) = \srule$ for each $1\leq i \leq n$ with
$\srule \in \cnf(\T)$ of the form $\bigwedge_{i=1}^n \fcu'_i\to \fml'$; and
\item $\fml = \bigvee_{i=1}^n\fmm_i\vee\fml'\s$ with $\s$ a MGU of $\fcu_i, \fcu'_i$ for each $1\leq i \leq n$.
\end{itemize}
Consider $\fct\in\fml$. To show property $1$ we need to find
a proof of $\Gamma_{\theta}(\fct\nu)$ in $\prog^{\chi}\cup\abox_0$. 
If $\fct\in\fmm_i$ then by i.h. we can find such a proof.
If  $\fct\in\fml'\s$ then it must be $\fct = \fct'\s$ for some $\fct'\in\fml'$ and, 
by definition of $\prog^{\chi}$, $\cnf(\T)$, and $\Gamma_{\theta}$, $\srule \in \cnf(\T)$ implies 
$\bigwedge_{i=1}^n \fcu'_i\to \Gamma_{\theta}(\fct')\in\prog^{\chi}$.
Since $\s$ is a MGU of $\fcu_i, \fcu'_i$ (with $\fcu_i = \fcu'_i\s$) for each $1\leq i \leq n$, 
$\s\nu$ must be a MGU of $\fcu_i\nu, \fcu'_i$ (with $\fcu_i\nu = \fcu'_i\s\nu$) for each $1\leq i \leq n$;
furthermore, since $\fcu'_i$ is necessarily function-free, 
it is $\Gamma_{\theta}(\fcu'_i\s\nu) = \fcu'_i\s\nu$, 
and thus $\Gamma_{\theta}(\fcu_i\nu) = \fcu'_i\s\nu$ and
$\s\nu$ is also a MGU of $\Gamma_{\theta}(\fcu_i\nu), \fcu'_i$ for each $1\leq i \leq n$.
By i.h. we have a proof in $\prog^{\chi}\cup\abox_0$ of $\Gamma_{\theta}(\fcu_i\nu)$ 
for each $1\leq i \leq n$;
it is easy to see that 
$\Gamma_{\theta}(\fct')\s\nu = \Gamma_{\theta}(\fct'\s\nu) = \Gamma_{\theta}(\fct\nu)$,
so combining these proofs with rule 
$\bigwedge_{i=1}^n \fcu'_i\to \Gamma_{\theta}(\fct')$ yields a proof of
$\Gamma_{\theta}(\fct\nu)$ in $\prog^{\chi}\cup\abox_0$.

Now consider $\ax\in\T$ such that $\rho$ mentions some $\srule'\in\cnf(\ax)$.
To show property 2 we need to find $\fct\in\fml\cup\set{\bot}$ and a proof of 
$\Gamma_{\theta}(\fct\nu)$ that mentions some rule in $\Xi^{\chi}(\rrule)$.
Assume first that $\srule' = \srule$.
If $\fml'=\emptyset$ then it must be 
$\cnf(\ax) = \set{\bigwedge_{i=1}^n \fcu'_i\to \bot}\subseteq\Xi^{\chi}(\rrule)$ and, as before,
we can combine this rule with proofs in $\prog^{\chi}\cup\abox_0$ of the $\Gamma_{\theta}(\fcu_i\nu)$ 
to obtain a proof of $\bot$ in $\prog^{\chi}\cup\abox_0$.
If $\fml'\neq \emptyset$ then it must be 
$\mset{\bigwedge_{i=1}^n \fcu'_i\to \Gamma_{\theta}(\fct')}{\fct'\in\fml'}\subseteq\Xi^{\chi}(\ax)$.
Since $\fml = \bigvee_{i=1}^n\fmm_i\vee\fml'\s$, for each $\fct'\in\fml'$ it is $\fct'\s\in\fml$ 
and, as we just saw, we can construct a proof of $\Gamma_{\theta}(\fct'\s\nu)$ that mentions
$\bigwedge_{i=1}^n \fcu'_i\to \Gamma_{\theta}(\fct')$.
%
Finally, assume that $\srule' \neq \srule$.
Then there must be some $i\in\set{1,\dots,n}$ such that $\srule'$ is mentioned by the proof 
$\rho_i$ of $\fcu_i\vee\fmm_i$ that is embedded in $\rho$.
Since $\rho_i$ is of depth $<d$, by i.h. there must be $\fcu''\in\fcu_i\vee\fmm_i$
and a proof $\rho''$ of $\Gamma_{\theta}(\fcu''\nu)$ in $\prog^{\chi}\cup\abox_0$ that mentions
some rule in $\Xi^{\chi}(\ax)$. If $\fcu'' \in \fmm_i$ then $\fcu''\in \fml$ already; if $\fcu'' = \fcu_i$ 
then, as before, for any $\fct\in\fml$ we can construct a proof of $\Gamma_{\theta}(\fct\nu)$ 
in $\prog^{\chi}\cup\abox_0$ such that $\rho''$ is embedded in it.
\qedhere
\end{itemize}
\end{proof}


\begin{restatable}{proposition}{generalConsequencePreservation}
\label{prop:generalConsequencePreservation}
Let $\rrule=\fml(\x)\to\fmm(\x)$ with $\fml$ a conjunction and $\fmm$ a disjunction
of atoms. Let $\upchi = \langle \theta, \abox_0, \arel\rangle$ be  a \bla{} satisfying
$\{\bot\} \subseteq \arel$ and such that for every substitution $\s$ mapping all variables in 
$r$ to pairwise distinct constants not in $\T$ there exists a mapping $\nu_{\s}$ with 
$\fml\s\nu_{\s}\subseteq\abox_0$ and $\fmm\s\nu_{\s}\subseteq\arel$. 
Then
\begin{enumerate}
\item $\T\models \rrule$ iff $\M^{\chi}\models \rrule$;

\item if $\T'\subseteq \T$ is a justification for $\rrule$ in $\T$ then $\T'\subseteq\M^{\chi}$;

\item $\T\setminus \M^{\chi}\models \rrule$ iff $\emptyset\models \rrule$.
\end{enumerate}
\end{restatable}
\begin{proof}\ 

\begin{enumerate}
\item By monotonicity, it is immediate that $\T\models \rrule$ if $\M^{\chi}\models \rrule$.

%Let us prove that $\T\models r$ implies $\M^{\chi}\models r$. 
Suppose $\T\models r$ and 
let $\s$ be a substitution mapping all variables in $\rrule$ to fresh, pairwise 
distinct constants. Then we have that $\T\cup\mset{\fct\s}{\fct\in\fml}\models \fmm\s$, which
implies $\cnf(\T)\cup\mset{\fct\s}{\fct\in\fml}\models \fmm\s$ and by
completeness of hyperresolution $\cnf(\T)\cup\mset{\fct\s}{\fct\in\fml}\vdash\fmm'$
for some $\fmm'\subseteq\fmm\s$.
Since $\mset{\fct\s\nu_{\s}}{\fct\in\fml}\subseteq \abox_0$, 
by Lemma \ref{lemma:HyperresolutionProofsToDatalogProofs} we have that for each 
$\srule\in\T$ such that some $\prule\in\cnf(\srule)$ supports 
$\fmm'$ in $\cnf(\T)\cup\mset{\fct\s}{\fct\in\fml}$
there exists $\fct\in\Gamma_{\theta}(\fmm'\nu_{\s})\cup\set{\bot}$
that is supported in $\prog^{\chi}\cup\abox_0$ by some rule from $\Xi^{\chi}(\srule)$.
By assumption, $\bot\in\arel$; also, $\fmm'$ is function-free so 
$\Gamma_{\theta}(\fmm'\nu_{\s}) = \fmm'\nu_{\s}$, and hence,
since $\fmm\s\nu_{\s}\subseteq \arel$, 
we have that $\fmm'\nu_{\s}\subseteq \arel$ and also $\fct\in\arel$.
In either case we have $\srule\in\M^{\chi}$ and consequently $\M^{\chi}\models \rrule$.

\item Let $\T'\subseteq \T$ be a justification for $\rrule$ in $\T$. 
As before, if $\s$ is a ground substitution for $\rrule$ mapping variables in $\rrule$ to fresh, 
pairwise distinct constants, then 
$\cnf(\T')\cup\mset{\fct\s}{\fct\in\fml}\vdash\fmm'$ 
for some $\fmm'\subseteq\fmm\s$.
In fact, by minimality of justifications, for each $\srule\in\T'$ some 
$\prule\in \cnf(\srule)$ must be in the support of 
some $\fmm'\subseteq \fmm\s$
in $\cnf(\T')\cup \mset{\fct\s}{\fct\in\fml}$.
As before, by Lemma \ref{lemma:HyperresolutionProofsToDatalogProofs}, this implies
$\srule\in\M^{\chi}$.

\item By monotonicity, it is immediate that 
$\T\setminus \M^{\chi}\models r$ if $\emptyset\models r$.

%Let us now prove that $\T\setminus\M^{\chi}\models r$ implies $\emptyset\models r$.
Suppose $\T\setminus\M^{\chi}\models \rrule$ and let $\T'$ be a justification for $\rrule$ in 
$\T\setminus\M^{\chi}$. Then $\T'$ is also a justification for $\rrule$ in $\T$, and, as we just 
proved, $\T'\subseteq \M^{\chi}$. This implies that $\T' = \emptyset$, and therefore
$\emptyset\models \rrule$.\qedhere
\end{enumerate}
\end{proof}


\begin{proposition}
\label{prop:cqPreservation}
Let $\ax = \fml(\x)\to \exists \y.[\bigvee_{i=1}^n \fmm_i(\x,\y)]$ be a rule. 
Let $\upchi = \langle \theta, \abox_0, \arel \rangle$ be a \bla{} satisfying
\begin{itemize}
\item $\{\bot\}\subseteq \arel$ and also $\fmm\s\subseteq \arel$ for every substitution 
$\s$ mapping all variables in $\rrule$ to constants in $\upchi$;
\item for every substitution $\s$ mapping all variables in $\rrule$ to pairwise distinct constants not in $\T$ there exists 
a mapping $\nu_{\s}$ such that $\fml\s\nu_{\s}\subseteq\abox_0$. 
\end{itemize}
Then
\begin{enumerate}
\item $\T\models \ax$ iff $\M^{\chi}\models \ax$;
\item if $\T'\subseteq \T$ is a justification for $\ax$ in $\T$ then $\T'\subseteq\M^{\chi}$;
\item $\T\setminus\M^{\chi}\models \ax$ iff $\emptyset\models \ax$.
\end{enumerate}
\end{proposition}


\begin{proof}\
\begin{enumerate}
\item By monotonicity, it is immediate that $\T\cup\abox\models \ax$ if $\M^{\chi}\cup\abox\models \ax$.

Let $\roleQ$ be a fresh predicate and $\T_{\fmm\to\roleQ} = \mset{\fmm_i(\x,\y) \to \roleQ(\x)}{1\leq i \leq n}$.
Then 
\[\T\models \ax \text{\  iff \ } 
\T\cup \T_{\fmm\to\roleQ} \models \fml(\x)\to\roleQ(\x)
\text{\ \ \ \ and \ \ \ \ }
\M^{\chi}\models \ax \text{\  iff \ } 
\M^{\chi}\cup\T_{\fmm\to\roleQ}\models \fml(\x)\to\roleQ(\x)\]

Consider $\T' = \T\cup\T_{\fmm\to\roleQ} $ and $\S' = \S\cup\{\roleQ\}$.
Clearly, $\T'$ has the exact same existentially quantified variables as $\T$. Therefore
$\upchi' = \langle \theta, \abox_0, \arel'\rangle$ with 
\[\arel' = \mset{\roleQ(\x)\s}{\s \text{ is a substitution mapping all variables in } \x \text{ to constants in }\upchi}\]
is a \bla{} for $\T'$ and $\S'$
and by Proposition \ref{prop:generalConsequencePreservation} we have that $\T'\models \fml(\x)\to\roleQ(\x)$ iff
$\M^{\chi'}\models \fml(\x)\to\roleQ(\x)$. 

If we show that $\M^{\chi'}\setminus\T_{\fmm\to\roleQ}  \subseteq \M^{\chi}$
then, by monotonicity, we will be able to conclude that 
\[\M^{\chi'}\models \fml(\x)\to\roleQ(\x) \text{ implies }
\M^{\chi}\cup\T_{\fmm\to\roleQ} \models \fml(\x)\to\roleQ(\x)\]
and thus that $\T\models \ax$ implies $\M^{\chi}\models \ax$.

Let $\srule\in \M^{\chi'}\setminus\T_{\fmm\to\roleQ} $.
Some $\prule\in\Xi^{\chi'}(\srule) = \Xi^{\chi}(\srule)$
must be in the support of some $\roleQ(\x)\s\in\arel'$ in $\prog^{\chi'}\cup\abox_0$.
In particular, $\prule$ must be mentioned in some proof $\rho = (T,\lambda)$ of $\roleQ(\x)\s$ 
in $\prog^{\chi'}\cup\abox_0$.
Let $v$ be the root of $T$ and $w_1,\dots,w_m$ its children nodes, there must be some 
$\bigwedge_{j=1}^m\fct_j(\x,\y)\to \roleQ(\x)\in\T_{\fmm\to\roleQ} $ 
and a MGU $\s'$ of $\fct_j, \lambda(w_j)$ for each $1\leq j \leq m$.
Since $\srule\notin \T_{\fmm\to\roleQ}$, 
there must exist $j\in\set{1,\dots,m}$ such that $\prule$ is mentioned in the proof $\rho_j$ 
of $\lambda(w_j)$ in $\prog^{\chi'}\cup\abox_0$ that is embedded in $\rho$.
Furthermore, since $\roleQ$ does not occur in the body of any rule in 
$\prog^{\chi'} = \prog^{\chi}\cup\T_{\fmm\to\roleQ}$, all rules 
mentioned in $\rho_j$ must be in $\prog^{\chi}$ and thus $\rho_j$ is a proof of $\fct$ in
$\prog^{\chi}\cup\abox_0$. Since by assumption $\lambda(w_j) = \fct_j\s'\in\arel$, this 
implies $\srule\in\M^{\chi}$.

\item 
Let $\T''\subseteq \T$ be a justification for $\ax$ in $\T$.
As before,
$\T''\models \ax \text{\ \  implies\ \ } 
\T''\cup\T_{\fmm\to\roleQ}\models \fml(\x)\to\roleQ(\x)$
and in particular for any substitution $\s$ mapping variables in $\rrule$ to pairwise distinct constants
we have 
$\T''\cup\T_{\fmm\to\roleQ} \cup \mset{\fct\s}{\fct\in\fml}\models \roleQ(\x)\s$
and therefore
$\cnf(\T''\cup\T_{\fmm\to\roleQ}) \cup \mset{\fct\s}{\fct\in\fml}\vdash \roleQ(\x)\s$.
%
%It is easy to see that for each $\srule$ in the support of $\roleQ(\x)\s$ in 
%$\cnf(\T''\cup\T_{\fmm\to\roleQ}) \cup \mset{\fct\s}{\fct\in\fml}$,
%$\srule$ must also be in the support of $\roleQ(\x)\s\nu_{\s}$ in 
%$\cnf(\T'\cup\T_{\fmm\to\roleQ}) \cup \mset{\fct\s\nu_{\s}}{\fct\in\fml}$.
%%Given a proof $\rho$ of $\roleQ(\x)\s$ in 
%%$\cnf(\T''\cup\T_{\fmm\to\roleQ}) \cup \mset{\fct\s}{\fct\in\fml}$, 
%%we can easily obtain a proof $\rho'$ of $\roleQ(\x)\s\nu_{\s}$ in 
%%$\cnf(\T'\cup\T_{\fmm\to\roleQ}) \cup \mset{\fct\s\nu_{\s}}{\fct\in\fml}$ with the same support as $\rho$.
By minimality of justifications, for each $\srule\in\T''$ there must be some $\prule\in\cnf(\srule)$
in the support of $\roleQ(\x)\s$ in $\cnf(\T''\cup\T_{\fmm\to\roleQ})\cup\mset{\fct\s}{\fct\in\fml} $.
It is easy to see that $\prule$ must also be in the support of $\roleQ(\x)\s\nu_{\s}$ in 
$\cnf(\T''\cup\T_{\fmm\to\roleQ}) \cup \mset{\fct\s\nu_{\s}}{\fct\in\fml}$.
Since $\roleQ(\x)\s\nu_{\s}\subseteq \arel'$
and $\mset{\fct\s\nu_{\s}}{\fct\in\fml}\subseteq \abox_0$,
by Lemma \ref{lemma:HyperresolutionProofsToDatalogProofs}
we have  that $\srule\in\M^{\chi'}$. In particular, since $\srule\in\T''\subseteq \T$, 
it must be 
$\srule\in\M^{\chi'} \setminus \T_{\fmm\to\roleQ} \subseteq \M^{\chi}$.

\item Again by monotonicity, it is immediate that  $\T\setminus\M^{\chi}\cup\abox\models \rrule$ if $\abox\models \rrule$.
By a similar argument to the one given in Proposition \ref{prop:generalConsequencePreservation}, it follows from 2 
that any justification for $\ax$ in $\T\setminus\M^{\chi}$ must be empty and therefore 
if $\T\setminus \M^{\chi}\models \ax$ then $\emptyset\models \ax$.\qedhere
\end{enumerate}
\end{proof}


\datalogAtomicSubsumptionModules*
\begin{proof}
  Consider an arbitrary rule of the form $\A(\x)\to \B(\x)$ with
  $\A,\B\in\S$ and $\A\ne\B$ (if $\A=\B$ the rule is tautological).
  Since $\x$ is implicitly universally quantified, we can assume
  w.l.o.g.  that $\x = (x_1,\dots,x_n)$ with $x_1, \dots,x_n$ pairwise
  distinct. %(and $n=\arity(\A)$).
  Let $\s$ be a substitution mapping $x_1,\dots,x_n$, respectively, to
  $c_1,\dots,c_n$, pairwise distinct constants not in $\T$.  Now
  consider a mapping $\nu_{\s}$ such that $c_i\nu_{\s} =
  c^i_{\A}$. This mapping is well-defined because $c_1,\dots,c_n$ are
  pairwise distinct. By definition of $\upchiimp$, we have
  $\A(\x)\s\nu_{\s}\in\abox_0^{\mathsf{i}}$ and
  $\B(\x)\s\nu_{\s}\in\arel^{\mathsf{i}}$, and therefore, by
  Proposition \ref{prop:generalConsequencePreservation}, we have
  $\T\models \A(\x)\to\B(\x)$ iff $\M^{\chiimp}\models
  \A(\x)\to\B(\x)$.
\end{proof}

\disjDatalogModules*
\begin{proof}
  By Proposition \ref{prop:fq-insep-charact-rules} it suffices to show
  that for any datalog rule $r=\fml\to\fmm$ we have $\T\models r$ iff
  $\M^{\chidd}\models r$.  Let $\s$ be a substitution mapping all
  variables in $r$ to pairwise distinct constants not in $\T$.
  Consider a mapping $\nu^*$ such that $x\s\nu^* = *$ for each
  $x\in\x$. Clearly $\fml\s\nu^*\subseteq\abox_0^{\mathsf{f}}$ and
  $\fmm\s\nu^*\subseteq\arel^{\mathsf{f}}$, and therefore, by
  Proposition \ref{prop:generalConsequencePreservation}, we have
  $\T\models r$ iff $\M^{\chidd}\models r$.
\end{proof}

\datalogCQModules*
\begin{proof}
  By Proposition \ref{prop:fq-insep-charact-rules} it suffices to show
  that for any rule $r=\fml\to\fmm$ we have $\T\models r$ iff
  $\M^{\chicq}\models r$.  Let $\s$ be a substitution mapping all
  variables in $r$ to pairwise distinct constants not in $\T$.
  Given a mapping $\nu^*$ such that $x\s\nu^* = *$ for each
  $x\in\x$ it is clear that
  $\fml\s\nu^*\subseteq\abox_0^{\mathsf{q}}$.  It is also immediate
  that $\fmm\s'\subseteq\arel^{\mathsf{q}}$ for every substitution
  $\s'$ mapping all variables in $r$ to constants in $\upchicq$.
  Therefore, by Proposition \ref{prop:cqPreservation}, we have 
  $\T\models r$ iff $\M^{\chicq}\models r$.
\end{proof}

\datalogClassifModules*
\begin{proof}
Analogous to the proof of Theorem \ref{thm:datalogAtomicSubsumptionModules}.
\end{proof}










\section{Model Inseparability}

Given an ABox $\abox$ and a datalog program $\prog$, let $\prog(\abox)$ denote the
materialisation of $\prog\cup\abox$. 
Furthermore, given a \bla{} $\upchi$, let $\rel(\upchi)$
denote the support of $\upchi$.

\datalogSemModules*
\begin{proof}
Let $\calI$ be a model of $\M^{\chistar}$. We assume w.l.o.g that $\calI$ is defined over all of $\sig{\T}$. 
Consider the interpretation $\calJ$ over $\sig{\T}$ such that  $\Delta^{\calJ} = \Delta^{\calI}$ and 
\[\begin{array}{rcl}
 \A^{\calJ} & = & 
\left\lbrace\begin{array}{lll}
\A^{\calI} &  \text{ if } \A\in(\S\cup\sig{ 
\rel(\upchi)})\setminus \set{\bot}\\
\D^{\arity(\A)} & \text{ if } \A\in\sig{\prog^{\chistar}(\abox_0^{\mathsf{m}})}\setminus (\S\cup\sig{\rel(\upchi)})\\
\emptyset & \text{ otherwise}
\end{array}
\right.
\end{array}\]
%Note that $\calJ$ is well defined 
%since $\sig{\rel(\upchi)}\setminus\set{\bot} \subseteq \sig{\M^{\chistar}}$.

%Consider $\ax : \bigwedge_{i=1}^n \fct_i(\x) \to \exists \y. [\bigvee_{j=1}^m  \fmm_j(\x,\y)] \in \T$.
%We will show that $\calJ\models \ax$.
Consider $\ax : \fml(\x) \to \exists \y. [\bigvee_{j=1}^m  \fmm_j(\x,\y)] \in \T$.
We will show that $\calJ\models \ax$.

Assume first $m=0$. Then $\Xi^{\chistar}(\ax) =\set{ \fml \to \bot}$.
If $r\in\M^{\chistar}$ then in particular $\sig{r}\subseteq \sig{\rel(\upchistar)}$, 
so $\calI$ and $\calJ$ agree over $\sig{\ax}$, and $\calJ\models \ax$.
If $r\notin\M^{\chistar}$ then, since $\bot\in\arel^{\chistar}$ and the only constant mentioned 
in $\prog^{\chistar}\cup\abox_0^{\mathsf{m}}$ is $*$, there must be $\fct\in\fml$ such that 
$\fct*\notin \prog^{\chistar}(\abox_0^{\mathsf{m}})$
(where, in an abuse of 
notation, $*$ denotes the substitution that maps all variables to $*$), and in particular 
$\sig{\fct}\not\subseteq \sig{\prog^{\chistar}(\abox_0^{\mathsf{m}})}$.
Since $\S\cup\sig{\rel(\upchistar)}\subseteq \sig{\prog^{\chistar}(\abox_0^{\mathsf{m}})}$,
this implies that for $\A\in\sig{\fct}$ it is $\A^{\calJ} = \emptyset$ and therefore
trivially $\calJ\models r$.

Assume now $m>0$ and let $\s$ be a substitution over all variables in $\ax$ such that 
%$\calJ\models\bigwedge_{\fct\in\fml} \fct\s$
$\calJ\models\fml\s$
(if no such substitution exists then 
trivially $\calJ\models \ax$). 
Since $\S\cup\sig{\rel(\upchi)}\subseteq \sig{\prog^{\chistar}(\abox_0^{\mathsf{m}})}$,
all predicates in $\fml$ must occur in $\prog^{\chistar}(\abox_0^{\mathsf{m}})$.
%The only constant mentioned in $\prog^{\chistar}\cup\abox_0^{\mathsf{m}}$ is $*$, so 
In particular it must be $\fct*\in\prog^{\chistar}(\abox_0^{\mathsf{m}})$ for every $\fct\in\fml$. 
This implies $\fcu*\in\prog^{\chistar}(\abox_0^{\mathsf{m}})$ for every $\fcu\in\bigcup_{j=1}^m\fmm_j$
and therefore for every predicate $\A$ in 
$\sig{\bigvee_{j=1}^m  \fmm_j}$ we have that either $\A^{\calJ} = \A^{\calI}$
or $\A^{\calJ} = \Delta^{\arity(\A)}$---in particular $\A^{\calI} \subseteq \A^{\calJ}$.
%
If $\A^{\calJ} = \Delta^{\arity(\A)}$ for every $\A\in\sig{\bigvee_{j=1}^m \fmm_j}$,
then it is immediate that $\calJ\models \ax$. Suppose there exists 
$\A\in\sig{\bigvee_{j=1}^m \fmm_j}$ such that $\A^{\calJ} \neq \Delta^{\arity(\A)}$.
Then $\A\in\S\cup\sig{\rel(\upchistar)}$.
%
If $\A\in\S$ then $\A(*,\dots,*)\in\arel^{\chistar}$.
Since $\A\in\sig{\bigvee_{j=1}^m\fmm_j}$ and 
$\fct*\in\prog^{\chistar}(\abox_0^{\mathsf{m}})$ for every $\fct\in\fml$, 
there is a proof $\rho_{\A,\ax}$ of $\A(*,\dots,*)$ in 
$\prog^{\chistar}\cup\abox_0^{\mathsf{m}}$ 
that mentions a rule in $\Xi^{\chistar}(\ax)$. Therefore $\ax\in\M^{\chistar}$.
%
If $\A\in\sig{\rel(\upchistar)}\setminus\S$ then some other $\fct'\in\arel^{\mathsf{m}}$
must be supported by a rule that has $\A$ in its signature.
More specifically, there must be a proof of $\fct'$ in $\prog^{\chistar}\cup\abox_0^{\mathsf{m}}$
that has a proof of $\A(*,\dots,*)$ as a subproof.
Replacing this subproof with $\rho_{\A,\ax}$ results in a proof of $\fct'$ in 
$\prog^{\chistar}\cup\abox_0^{\mathsf{m}}$ that mentions a rule in $\Xi^{\chistar}(\ax)$. 
Therefore in this case $\ax\in\M^{\chistar}$ too.
%
Now, since all rules in $\Xi^{\chistar}(\ax)$ have the same body as $\ax$, we have that 
$\sig{\fml}\subseteq \rel(\upchistar)\setminus\set{\bot}$ and therefore $\calI$ and $\calJ$ agree over $\sig{\fml}$. 
By assumption, $\calJ\models\fml\s$, so also $\calI\models\fml\s$;
furthermore $\calI\models\M^{\chistar}$ implies
$\calI\models\bigvee_{j=1}^m  \fmm_j\s$, which implies 
$\calJ\models\bigvee_{j=1}^m  \fmm_j\s$ because 
$\A^{\calI} \subseteq \A^{\calJ}$
for every predicate $\A\in\sig{\bigvee_{j=1}^m  \fmm_j}$. 
Since $\s$ is arbitrary, we conclude that $\calJ\models \ax$.
\end{proof}













\section{Depletingness and Preservation of Justifications}

\depletingness*
\begin{proof}
For $z\in\set{\mathsf{q},\mathsf{f},\mathsf{i},\mathsf{c}}$
the statement follows from Propositions \ref{prop:fq-insep-charact-rules}, 
\ref{prop:generalConsequencePreservation} and 
\ref{prop:cqPreservation} by the arguments already presented in the proofs for 
Theorems \ref{thm:datalogAtomicSubsumptionModules}, 
\ref{thm:disjDatalogModules} and \ref{thm:datalogCQModules}.

For $z=\mathsf{m}$, we will now show that $\T\setminus \M^{\chistar} \eqm \emptyset$.
Let $\calI$ be a model of $\emptyset$ and $\abox_{\calI}$ the ABox defined by $\calI$ over $\S$. 
Consider the datalog program $\prog = \bigcup_{\ax\in\T\setminus\M^{\chistar}}\Xi_{\chistar}(\ax),$
and the materialisation $\prog(\abox_{\calI})$ of $\prog$ w.r.t. $\abox_{\calI}$. 
We show that $\prog(\abox_{\calI})$ is a model of $\T\setminus\M$ that coincides with $\calI$ over $\S$. For this, it 
suffices to show the following two properties:
\begin{itemize}
\item All facts over $\S$ in $\prog(\abox_{\calI})$ must already be in $\abox_{\calI}$.\\
  Let $\fct\in\prog(\abox_{\calI})$ be a fact over $\S$. If
  $\fct\notin\abox_{\calI}$ then there must exist a proof $\rho$ of
  $\fct$ in $\prog\cup\abox_{\calI}$. Since $\abox_{\calI}$ only
  mentions predicates from $\S$ and $\prog\subseteq \prog^{\chistar}$,
  we can find a proof of $\fct*\in\arel^{\mathsf{m}}$ in
  $\prog^{\chistar}\cup\abox_0^{\mathsf{m}}$ that mentions the exact
  same rules as $\rho$.  Let $\rrule$ be a rule mentioned in $\rho$,
  there must exist $\srule\in\T\setminus\M^{\chistar}$ such that
  $\rrule\in\Xi^{\chistar}(\srule)$; however, because $\rrule$ is also
  mentioned in a proof of $\fct*$ in
  $\prog^{\chistar}\cup\abox_0^{\mathsf{m}}$, it must also be
  $\srule\in\M^{\chistar}$.  This is a contradiction, so
  $\fct\in\abox_{\calI}$.

\item $\bot\notin\prog(\abox_{\calI})$.\\
Suppose $\bot\in\prog(\abox_{\calI})$. Then there must be a proof $\rho$ of $\bot$ in 
$\prog\cup\abox_{\calI}$. Again, we can find a proof of $\bot$ in $\prog^{\chistar}\cup\abox_0^{\mathsf{m}}$ 
supported by the exact same rules as $\rho$. Following a similar argument as before, we 
conclude that $\bot\notin\prog(\abox_{\calI})$.\qedhere
\end{itemize}
%
\end{proof}

\presJustifications*
\begin{proof}
The claim follows from Propositions \ref{prop:fq-insep-charact-rules}, 
\ref{prop:generalConsequencePreservation} and \ref{prop:cqPreservation}
similarly to Theorems \ref{thm:datalogAtomicSubsumptionModules}, 
\ref{thm:disjDatalogModules} and \ref{thm:datalogCQModules}.
\end{proof}












\section{Module Containment}

\begin{definition}
\label{def:blaHomomorphism}
Let $\upchi = \langle \theta, \abox_0, \arel \rangle$ and $\upchi' = \langle \theta', \abox'_0, \arel' \rangle$ and let $\N$ and $\N'$ be the sets of constants mentioned in $\upchi$ and $\upchi'$, respectively.
A mapping $\mu:\N \to \N'$ is a \emph{homomorphism from $\upchi$ to $\upchi'$} if 
the following conditions hold:
\begin{inparaenum}[(i)]
\item $\theta' = \theta\mu$,
\item $\abox_0\mu \subseteq \abox'_0 $; and 
\item $\arel\mu \subseteq \arel'$. 
\end{inparaenum}
We write
$\upchi\hookrightarrow\upchi'$ if a homomorphism from $\upchi$ to $\upchi'$ exists.
\end{definition}

\begin{restatable}{theorem}{homorphismSubset} \label{thm:homomorphism-subset}
  If $\upchi,\upchi'$ are s.t. $\upchi\hookrightarrow\upchi'$, then
  $\M^{\chi} \subseteq \M^{\chi'}$.
\end{restatable}

%\homorphismSubset*
\begin{proof}
Suppose $\upchi = \langle \theta, \abox_0, \arel \rangle$ and $\upchi' = \langle\theta', \abox_0', \arel' \rangle$ 
with $\N$  and $\N'$ the sets of constants mentioned in $\upchi$  and $\upchi'$, respectively.
Let $\mu$ be a homomorphism from $\upchi$ to $\upchi'$ and let $\ax\in\M^{\chi}$. 
Some $\fct\in\arel$ %such that some proof of $\fct$ in $\prog^{\chi}\cup\abox_0$ is 
must be supported  in $\prog^{\chi}\cup\abox_0$ by a rule in $\Xi^{\chi}(\ax)$. 
Since, by assumption, $\fct\mu\in\arel'$, it suffices for us to show that
%there exists a proof of 
$\fct\mu$ is supported in $\prog^{\chi'}\cup\abox'_0$ by a rule from $\Xi^{\chi'}(\ax)$.

%There must exist $r\in \Xi^{\chi}(\ax)$ and $\fct\in\arel$ such that $r$ supports a proof 
%$\rho = (T, \lambda)$ of $\fct$ in $\prog^{\chi}\cup\abox_0$. Let us show that there must also exist a proof for $\fct\mu$ 
%in $\prog^{\chi'}\cup\abox_0'$ that is supported by some rule in $\Xi^{\chi'}(\ax)$.

To this end we will show that for any rule $\ax\in\T$ and any fact $\fct$ such that there exists a proof 
$\rho = (T,\lambda)$ of $\fct$ in $\prog^{\chi}\cup\abox_0$ mentioning $\srule\in\Xi^{\chi}(\ax)$, 
there exists a proof of $\fct\mu$ in $\prog^{\chi'}\cup\abox'_0$ mentioning $\Xi^{\chi'}(\ax)$.
We will reason by induction on the depth $d$ of $\rho$---which must be at least $1$ since by assumption it uses $\srule$.

\begin{itemize}
\item [$d=1$] \ \\
$\ax$ must be of the form $\bigwedge_{i=1^n}\fcu'_i(\x) \to \exists \y. [\bigvee_{j=1}^m \fmm_j(\x,\y)]$
so $\srule$ must be 
	\begin{itemize}
	\item $\bigwedge_{i=1^n}\fcu'_i \to \fct'\theta$ with $\fct'\in\fmm_j$ for some $1\leq j \leq m$ if $m>0$ \\
          The $\lambda$-images of the leaves of $T$ must be
          $\fcu_1,\dots,\fcu_n \in \abox_0$ such that there exists a
          MGU $\s$ of $\fcu_i, \fcu'_i$ for every $1\leq i \leq n$
          satisfying $\fct = \fct'\theta\s$.
	%
	By assumption, we have $\fcu_i\mu\in\abox'_0$ for every $1\leq i \leq n$, and also
	$\srule' = \bigwedge_{i=1^n}\fcu'_i \to \fct'\theta' \in \Xi^{\chi'}(\ax)$ where $\theta'=\theta\mu$.
    %
	Consider $\s' = \s\mu$. 
	It is easy to see that $\mu\s\mu = \s\mu$ since the domain of $\s$ is disjoint with both the domain and 
	the range of $\mu$. Therefore $(\fcu_i\mu)\s\mu = \fcu_i\s\mu$ for every $1\leq i \leq n$.
	Furthermore, since $\s$ is a MGU of $\fcu_i, \fcu'_i$ for every $1\leq i \leq n$, we have that 
	$\s'$ is a MGU of $\fcu_i\mu, \fcu'_i$ for every $1\leq i \leq n$.
	Finally, since $\theta' = \theta\mu$, we have that 
	$\fct\mu = \fct'\theta\s\mu = \fct'\theta\mu\s\mu = \fct'\theta'\s'$ is a consequence of
	$\srule'$ and $\fcu_1\mu, \dots, \fcu_n\mu$, and hence we have a proof of 
	$\fct\mu$ in $\prog^{\chi'}\cup\abox_0'$ 
	supported by $\srule'\in\Xi^{\chi'}(\ax)$.
	
      \item  $\bigwedge_{i=1^n}\fcu'_i \to \bot$ if $m=0$ 	\\
	Then it must be $\fct = \bot$ and, as in the previous case,
        the $\lambda$-images of the leaves of $T$ must be
        $\fcu_1,\dots,\fcu_n \in \abox_0$ such that there exists a MGU
        $\s$ of $\fcu_i, \fcu'_i$ for every $1\leq i \leq n$.
	%
	Also, $\fcu_i\mu\in\abox'_0$ for every $1\leq i \leq n$, $\srule \in \Xi^{\chi'}(\ax)$, and 
	$\s' = \s\mu$ is a MGU of $\fcu_i\mu, \fcu'_i$ for every $1\leq i \leq n$, so we have a proof of 
	$\fct\mu = \bot$ in $\prog^{\chi'}\cup\abox_0'$ supported by $\srule\in\Xi^{\chi'}(\ax)$.
	\end{itemize}

      \item [$d>1$]\ \\
	Let $v$ be the root of $T$, let $\fcu_1,\dots,\fcu_n \in
        \abox_0$ be the $\lambda$-images of the children of $v$ and
        let $\ax'\in\T$ be such that the $\lambda$-image of the
        edges connecting $v$ with its children in $T$ is a rule in
        $\Xi^{\chi}(\ax')$.  Either $\srule\in\Xi^{\chi}(r')$ or it is
        mentioned in some subproof of $\rho$.  Our induction
        hypothesis implies that for each each $\ax''\in\T$, if some
        $\fcu_i$ is supported in $\prog^{\chi}\cup\abox_0$ by a rule
        in $\Xi^{\chi}(\ax'')$, then also $\fcu_i\mu$ is supported in
        $\prog^{\chi'}\cup\abox'_0$ by some rule in
        $\Xi^{\chi'}(\ax'')$.  Therefore, in either case, following an
        argument similar to case $d=1$, we can construct a proof
        $\rho'$ of $\fct\mu$ in $\prog^{\chi'}\cup\abox'_0$ from a
        collection of proofs of $\fcu_1\mu, \dots, \fcu_n\mu$ in
        $\prog^{\chi'}\cup\abox'_0$ and a rule in $\Xi^{\chi'}(\ax')$
        in such a way that $\rho'$ mentions a rule in
        $\Xi^{\chi'}(\ax)$.\qedhere
\end{itemize}
\end{proof}

\hierarchy*
\begin{proof}
%
This follows immediately from Theorem \ref{thm:homomorphism-subset}.
%
%For $\M^{\chiimp}\subseteq \M^{\chidd}$, 
%consider the mapping $\mu: \N^{\chiimp} \to \N^{\chidd}$ 
%such that $\mu(c_{\A}) = *$ for every $\A\in\S$ and 
%$\mu(c_x) = c_x$ 
%for each existentially quantified variable $x$ in $\T$. 
%
%For $\M^{\chidd}\subseteq \M^{\chicq}$
%consider the mapping $\mu: \N^{\chidd} \to \N^{\chicq}$ 
%such that $\mu(*) = *$ and $\mu(c_x) = c_x$ 
%for each existentially quantified variable $x$ in $\T$. 
%
%For $\M^{\chicq}\subseteq \M^{\chistar}$
%consider the mapping $\mu: \N^{\chicq} \to \{*\}$.
%%\end{proof}
%%
%%\hierarchytwo*
%%\begin{proof}
%For $\M^{\chistar}\subseteq \M^{\chibot}$ consider 
%the identity mapping on $\set{*}$.
%
%For $\chiimp \subseteq \chicls$ consider 
%the identity mapping on $\mset{c_{\A}}{\A\in\S}$.
%
%For $\chicls \subseteq \chibot$ consider 
%the mapping $\mu :\N^{\chicls} \to \set{*}$.
\end{proof}





% \begin{proposition}
% For any $x\in\set{\mathsf{i},\mathsf{f},\mathsf{q},\mathsf{m},\mathsf{c}}$, $\M^{\chi_i}$ is not strongly depleting.
% \end{proposition}
% \begin{proof}
% By Proposition \ref{prop:module-hierarchy},
% $\M^{\chiimp} {\subseteq} \M^{\chidd} {\subseteq} \M^{\chicq} {\subseteq} \M^{\chistar}$
% and $\M^{\chiimp}{\subseteq}\M^{\chicls}{\subseteq}\M^{\chistar}$,
% and by Proposition \ref{prop:insep-rel-hierarchy},
% $\eqm \subseteq \eqq \subseteq \eqf \subseteq \eqi$
% and $\eqc \subseteq \eqi$. Therefore, it is enough for us to provide an example of $\T$ and $\S$ such that
% $\T\setminus\M^{\chistar}{\not\eqi_{\S\cup\sig{\M^{\chistar}}}} \emptyset$.
% %
% Consider $\S=\set{\A,\D}$ and $\T = \set{ \ax_{8}, \ax_{9}, \ax_{10}, \ax_{11}}$ with
% %
% \begin{align}
% &\A(x) \to \exists r. [\roleR(x,y) \wedge \B(y)]\tag{$\ax_8$}\\
% &\A(x) \to \exists r. [\roleR(x,y) \wedge \C(y)]\tag{$\ax_9$}\\
% &\roleR(x,y) \to \D(x)\tag{$\ax_{10}$}\\
% &\B(x) \to \C(x)\tag{$\ax_{11}$}
% \end{align}
% %
% % \[\begin{array}{l}
% % \ax_{8} = \A(x) \to \exists r. [\roleR(x,y) \wedge \B(y)]\\
% % \ax_{9} = \A(x) \to \exists r. [\roleR(x,y) \wedge \C(y)]\\
% % \ax_{10} = \roleR(x,y) \to \D(x)\\
% % \ax_{11} = \B(x) \to \C(x)  \end{array} \]
% %
% Then $\M^{\chistar} {=} \set{\ax_{8}, \ax_{9}, \ax_{10}}$, and we have 
% $\B,\C \in \sig{\M^{\chistar}}$ and $\T\models \B(x) \to \C(x)$,
% but $\set{\ax_{8}, \ax_{9}, \ax_{10}} \not\models \B(x) \to \C(x)$ and 
% $\T\setminus\set{\ax_{8}, \ax_{9}, \ax_{10}} \models \B(x) \to \C(x)$.
% \end{proof}

%\section{Complexity}
%\complexity*
%\begin{proof}
%Both $\Xi^{\chi}(\T)$ and the materialisation of $\T$ w.r.t. $\abox_0$ can clearly be computed in time 
%polynomial in the size of $\T$. Furthermore, $\rel(\arel,\Xi^{\chi}(\T), \abox_0)$ can be computed in 
%time polynomial in the size of $\T$ too, using the technique presented in \cite{DBLP:conf/aaai/ZhouNGH14}.
%\end{proof}


%
%
%\section{Proofs for Section \ref{sec:addit-prop-modul}}
%
%\textbf{Do we need this in the end?}
%\begin{restatable}{theorem}{rewBotModules} \label{thm:datalogBotModules-1}
%  $\M^{\chibot}\eqm\T$.
%% Furthermore,
%%  $\M^{\chibot}$ is strongly depleting and self-contained.
%\end{restatable}
%%\rewBotModules*
%  $\M^{\chibot}\eqm\T$. Furthermore, $\M^{\chibot}$ is strongly depleting and
%  self-contained.
%\begin{proof}
%Consider any model $\calI = (\Delta,\cdot^{\calI})$ of $\M^{\chibot}$ and the interpretation 
%$\calJ = (\Delta, \cdot^{\calJ})$ of $\T$ such that $\A^{\calJ} = \A^{\calI}$ if 
%$\A\in\S\cup\sig{\M^{\upchi_1}}$ and $\A^{\calJ} = \emptyset$ otherwise.
%
%Let $\ax = \fml(\x) \to \exists \y. [\bigvee_{i=1}^n \fmm_i(\x,\y)]$ be an rule $\T$.
%If $\ax\in \M^{\chibot}$ then clearly $\calJ\models \ax$ since 
%$\calI|_{\sig{\M^{\chibot}}} = \calJ|_{\sig{\M^{\chibot}}}$.
%
%If $\ax\in\T\setminus \M^{\chibot}$ then no rule in $\Xi^{\chibot}(\ax)$ supports a proof in
%$\prog^{\chibot}\cup\abox_0$ of a fact in $\arel$. Note that, since the only constant mentioned
%in $\prog^{\chibot}\cup\abox_0$ is $*$, the materialisation of $\prog^{\chibot}$ w.r.t. $\abox_0$ 
%is necessarily contained in $\arel$. Therefore $\ax\in\T\setminus \M^{\chibot}$ implies that 
%there must be some $\A\in\sig{\fml}$ such that $\A\notin\sig{\M^{\chibot}}\cup\S$. But then 
%$\A^{\calJ} = \emptyset$ and trivially $\calJ\models \ax$.
%Therefore $\T\eqm \M^{\chibot}$. 
%Furthermore, since $\calI|_{\S\cup\sig{\M^{\chibot}}} = \calJ|_{\S\cup\sig{\M^{\chibot}}}$
%we have that $\M^{\chibot}$ is self-contained.
%
%Following a similar argument we can see that for any interpretation $\calI = (\Delta,\cdot^{\calI})$ the
%interpreatation $\calJ = (\Delta, \cdot^{\calJ})$ of $\T\setminus\M^{\chibot}$ such that 
%$\A^{\calJ} = \A^{\calI}$ if $\A\in\S\cup\sig{\M^{\chibot}}$ and $\A^{\calJ} = \emptyset$ otherwise
%is a model of $\T\setminus\M^{\chibot}$, so $\M^{\chibot}$ is strongly depleting.
%\end{proof}





%%% Local Variables: 
%%% mode: latex
%%% TeX-master: "paper"
%%% End: 
